\documentclass[11pt,letterpaper]{article}
\usepackage[margin=1.5cm]{geometry}
\usepackage{fontspec}
\setmainfont{Arial}
\usepackage{polyglossia}
\setdefaultlanguage{spanish}
\usepackage{enumitem}
\usepackage{array,booktabs}
\usepackage{xcolor}
\usepackage{pgfgantt}
\usepackage[hidelinks]{hyperref}

\definecolor{wpone}{RGB}{52,152,219}
\definecolor{wptwo}{RGB}{46,204,113}
\definecolor{wpthree}{RGB}{155,89,182}
\definecolor{wpfour}{RGB}{241,196,15}
\definecolor{wpfive}{RGB}{231,76,60}

\setlength{\parskip}{0.35em}
\setlength{\parindent}{0.5pt}

\begin{document}

\section*{Memoria Técnica del Proyecto -- Becas Leonardo 2026}

\textbf{Título:} \textit{Desarrollo de Plataforma Computacional Integral de Bioimagen, Visión Artificial y Genómica para el Diagnóstico de la Disquinesia Ciliar Primaria} \\
\textbf{Área:} Ciencias de la Computación, Ciencia de Datos e Inteligencia Artificial \\

\hrule

\subsection*{1. Problema y justificación}

La \textbf{Disquinesia Ciliar Primaria (DCP)} es una enfermedad genética multisistémica caracterizada por la alteración funcional o ultraestructural de las cilias móviles. \textbf{Aunque históricamente fue catalogada como infrecuente (1:15.000 a 1:30.000), estudios genómicos recientes sitúan su prevalencia en al menos 1:7.500 recién nacidos} \cite{Robson2025,Collison2025}. Presenta herencia predominantemente autosómica recesiva y una marcada heterogeneidad genética vinculada a mutaciones en más de 60 genes axonémicos \cite{Collison2025,Wrona2025}. La falla en el batido ciliar impide el aclaramiento mucociliar eficaz, generando acumulación crónica de secreciones e infecciones respiratorias recurrentes que conducen de forma progresiva a daño pulmonar estructural irreversible (bronquiectasias) y deterioro de la función ventilatoria \cite{Robson2025,Shapiro2015,Wrona2025}.

El diagnóstico internacional estándar (guías ERS/ATS) requiere un algoritmo multimodal que articula índices clínicos, óxido nítrico nasal (nNO), videomicroscopía de alta velocidad (HSVM), microscopía electrónica de transmisión (TEM) y secuenciación genética \cite{Shapiro2015,Balinotti2026}. \textbf{En Argentina, la escasez de equipamiento especializado provoca un retraso diagnóstico crítico}: la cohorte del Laboratorio de Motilidad Ciliar del Hospital de Niños Dr. Ricardo Gutiérrez revela una mediana de edad al diagnóstico de 8,8 años (marcadamente superior a los 5,3 años europeos), dilatándose hasta los \textbf{10,5 años} en pacientes sin anomalías viscerales (\textit{situs solitus}) \cite{Collison2025,Balinotti2026}. A pesar de que los síntomas inician precozmente (rinorrea persistente en el 81\% y tos húmeda crónica en el 90\% antes de los 6 meses de vida), esta demora propicia tratamientos inadecuados (como asma refractaria en el 52\%) y daño bronquiectásico irreversible instalado al momento del diagnóstico en el 49\% de los pacientes pediátricos \cite{ComiteNeumonologia2020,Balinotti2026}.

Frente a la complejidad de las guías internacionales, el \textbf{Hospital Gutiérrez implementó y validó la primera estrategia diagnóstica combinada adaptada al sistema de salud público argentino} \cite{Balinotti2026}. Dicho protocolo articula cribado clínico (ATS-CSQ y PICADAR), medición de nNO y HSVM a partir de biopsias de cornete nasal, reservando la TEM y la secuenciación genética para casos no concluyentes o según disponibilidad de recursos. Esta estrategia optimizó los recursos hospitalarios al clasificar con alta probabilidad o certeza al 47\% de los derivados y descartar con seguridad la enfermedad en el 49\% de los casos sospechosos \cite{Balinotti2026}. No obstante, la batería presenta una limitación operativa determinante: la medición de nNO por quimioluminiscencia requiere maniobras de cooperación activa confiables solo en niños $\ge$5 años \cite{Shapiro2015,Wrona2025,Balinotti2026}. En lactantes y preescolares ---franja donde el diagnóstico precoz previene secuelas crónicas--- el nNO resulta inviable, haciendo que la confirmación clínica descanse ineludiblemente en la HSVM, la TEM y el análisis genético \cite{ComiteNeumonologia2020,Balinotti2026}.

\textbf{El presente proyecto aborda estos cuellos de botella implementando una plataforma computacional integral de bioimagen, visión artificial y genómica \textit{in-silico} adaptada al hospital público}. La incorporación de algoritmos de visión por computadora para HSVM (extracción de frecuencia [CBF] mediante FFT y patrones de batido [CBP] por flujo óptico) y de aprendizaje automático para TEM (\textit{convpaint}) sustituirá la inspección visual manual y subjetiva por métricas reproducibles e independientes del operador \cite{Bricmont2021,Wrona2025}. Complementariamente, un flujo bioinformático de bajo costo facilitará la priorización de variantes en >60 genes de DCP a partir de datos abiertos \cite{Paff2021,Collison2025}. El impacto del proyecto reside en transformar radicalmente la capacidad diagnóstica en la primera infancia: al brindar herramientas objetivas, rápidas y accesibles, se facilitará el diagnóstico temprano en niños menores de 5 años, evitando la progresión hacia bronquiectasias irreversibles y mejorando drásticamente el pronóstico y la calidad de vida pediátrica.

\medskip\hrule\medskip


\subsection*{2. Estado del arte, innovación e hipótesis}
Actualmente, el hospital cuenta únicamente con un microscopio óptico estándar de campo claro antiguo acoplado a una cámara no apta que adquiere a tasas convencionales (≤30 fps), provocando un submuestreo severo (\textit{aliasing}) frente a cilias que baten a 10–15 Hz \cite{Bricmont2021}. \textbf{Más crítico aún es la ausencia total de control ambiental}: a temperatura ambiente no controlada (20–22 °C), la frecuencia de batido ciliar se deprime sustancialmente (la CBF es de 6,3–9,0 Hz a 32 °C y de 10–15 Hz a 37 °C) y la muestra se deseca rápidamente en el portaobjetos, induciendo discinesias secundarias y falsos positivos \cite{Bricmont2021,Wrona2025}. \textbf{Finalmente, la evaluación clínica actual es visual y manual, lo que genera una alta subjetividad y dependencia del operador} \cite{Bricmont2021,Balinotti2026}.

Para superar esta barrera, este proyecto establece una sinergia estratégica orientada a la ingeniería, la computación científica y la capacitación médica:
\begin{itemize}[leftmargin=*]
\item \textbf{Sede de Ejecución y Cómputo (FCEN-UBA):} La Facultad de Ciencias Exactas y Naturales de la UBA lidera el desarrollo de los algoritmos de visión artificial, los flujos bioinformáticos \textit{in-silico} de bajo costo, el diseño y fabricación de la cámara ambiental termostatizada y la dirección de estudiantes de grado universitarios.
\item \textbf{Colaboración Hospitalaria (Hospital Gutiérrez):} Médicos especialistas de los Servicios de Neumonología y Patología del hospital colaboran en la definición de requisitos clínicos, pruebas del equipamiento óptico en banco de pruebas y reciben la capacitación metodológica para operar la plataforma.

\end{itemize}

El estado del arte demuestra que la frecuencia de batido ciliar (CBF) no debe interpretarse de forma aislada del patrón de batido (CBP), y que ni la TEM ni la genética son universalmente concluyentes si se usan de manera independiente \cite{Shapiro2015,Bricmont2021,Wrona2025}. Por lo tanto, la oportunidad de innovación no está en repetir pruebas existentes, sino en integrar adquisición estandarizada, análisis automático multimodal y reporte reproducible.

\textbf{Novedad del proyecto:} integrar, en un flujo único, (a) instrumentación térmica in situ para calidad del dato, (b) visión artificial para HSVM, (c) análisis cuantitativo de TEM y (d) bioinformática in-silico de bajo costo, con transferencia efectiva al hospital público.

\textbf{Hipótesis:} una plataforma computacional multimodal, con adquisición controlada y análisis automatizado, reducirá variabilidad inter-operador, mejorará trazabilidad diagnóstica y aumentará factibilidad de diagnóstico temprano en contexto pediátrico público \cite{Robson2025,Collison2025,Balinotti2026}.

\subsection*{3. Objetivos}

\textbf{Objetivo general:} desarrollar e implementar una plataforma de bioimagen, visión artificial y genómica para soporte diagnóstico de DCP, con transferencia tecnológica al Hospital Gutiérrez.

\textbf{Objetivos específicos:}
\begin{enumerate}[leftmargin=*]
\item Implementar estación de adquisición con control térmico/hídrico y cámara de alta velocidad (WP1).
\item Desarrollar pipeline de HSVM para CBF/CBP cuantitativos y reporte asistido (WP2).
\item Automatizar cuantificación ultraestructural en TEM (WP3).
\item Implementar pipeline bioinformático in-silico para priorización de variantes y reporte molecular (WP4).
\item Capacitar equipos clínicos y liberar herramientas en acceso abierto (WP5).
\end{enumerate}

\subsection*{4. Metodología por paquetes de trabajo (18 meses)}

\textbf{WP1 -- Instrumentación y adquisición (M1--M6)}
\begin{itemize}[leftmargin=*]
\item \textbf{WP1.1} Diseño y validación de cámara ambiental termostatizada para platina (37\,$^{\circ}$C, humedad alta).
\item \textbf{WP1.2} Integración de cámara de alta velocidad y configuración óptica para adquisición estandarizada.
\item \textbf{WP1.3} Implementación de flujo de captura en software abierto (Micro-Manager/Python).
\item \textbf{Entregable:} estación HSVM operativa y protocolo de adquisición reproducible.
\end{itemize}

\textbf{WP2 -- Visión artificial para HSVM (M4--M12)}
\begin{itemize}[leftmargin=*]
\item \textbf{WP2.1} Preprocesamiento y segmentación de epitelio ciliado activo.
\item \textbf{WP2.2} Estimación de CBF por FFT y mapas espaciales de frecuencia.
\item \textbf{WP2.3} Cuantificación de CBP por flujo óptico y kymographs; clasificación de disquinesias.
\item \textbf{Entregable:} módulo de análisis HSVM con salida cuantitativa para uso médico.
\end{itemize}

\textbf{WP3 -- Análisis de bioimágenes TEM (M6--M14)}
\begin{itemize}[leftmargin=*]
\item \textbf{WP3.1} Curaduría y estandarización de micrografías TEM.
\item \textbf{WP3.2} Segmentación axonémica semiautomática (convpaint + ML).
\item \textbf{WP3.3} Cuantificación de defectos ODA/IDA y geometría axonemal 9+2.
\item \textbf{Entregable:} reporte TEM cuantitativo reproducible.
\end{itemize}

\textbf{WP4 -- Bioinformática in-silico (M8--M16)}
\begin{itemize}[leftmargin=*]
\item \textbf{WP4.1} Curaduría de variantes en >50 genes de DCP desde repositorios abiertos.
\item \textbf{WP4.2} Pipeline de anotación/priorización (VEP, SnpEff, CADD, REVEL, AlphaMissense).
\item \textbf{WP4.3} Integración genotipo-fenotipo con resultados HSVM/TEM.
\item \textbf{Entregable:} pipeline reproducible y catálogo regional de variantes.
\end{itemize}

\textbf{WP5 -- Transferencia, capacitación y ciencia abierta (M12--M18)}
\begin{itemize}[leftmargin=*]
\item \textbf{WP5.1} Talleres para neumonología y patología en uso instrumental y lectura de reportes.
\item \textbf{WP5.2} Transferencia de SOPs y manuales de operación.
\item \textbf{WP5.3} Publicación del software y documentación en GitHub; informe final BBVA.
\item \textbf{Entregable:} adopción operativa hospitalaria y suite abierta publicada.
\end{itemize}

\subsection*{5. Cronograma (códigos trazables WPx.y)}

{\small
\begin{center}
\begin{ganttchart}[
    x unit=0.46cm,
    y unit chart=0.5cm,
    y unit title=0.55cm,
    vgrid,
    hgrid,
    bar height=0.42,
    bar label font=\scriptsize,
    group label font=\bfseries\scriptsize,
    group peaks height=0.2
]{1}{18}
\gantttitle{Cronograma de ejecución (18 meses)}{18} \\
\gantttitlelist{1,...,18}{1} \\
\ganttgroup[group/.append style={fill=wpone, draw=wpone!70!black}]{WP1}{1}{6} \\
\ganttbar[bar/.append style={fill=wpone!65, draw=wpone!70!black}]{WP1.1}{1}{4} \\
\ganttbar[bar/.append style={fill=wpone!65, draw=wpone!70!black}]{WP1.2}{2}{5} \\
\ganttbar[bar/.append style={fill=wpone!65, draw=wpone!70!black}]{WP1.3}{4}{6} \\
\ganttgroup[group/.append style={fill=wptwo, draw=wptwo!70!black}]{WP2}{4}{12} \\
\ganttbar[bar/.append style={fill=wptwo!65, draw=wptwo!70!black}]{WP2.1}{4}{7} \\
\ganttbar[bar/.append style={fill=wptwo!65, draw=wptwo!70!black}]{WP2.2}{6}{10} \\
\ganttbar[bar/.append style={fill=wptwo!65, draw=wptwo!70!black}]{WP2.3}{8}{12} \\
\ganttgroup[group/.append style={fill=wpthree, draw=wpthree!70!black}]{WP3}{6}{14} \\
\ganttbar[bar/.append style={fill=wpthree!65, draw=wpthree!70!black}]{WP3.1}{6}{8} \\
\ganttbar[bar/.append style={fill=wpthree!65, draw=wpthree!70!black}]{WP3.2}{8}{11} \\
\ganttbar[bar/.append style={fill=wpthree!65, draw=wpthree!70!black}]{WP3.3}{11}{14} \\
\ganttgroup[group/.append style={fill=wpfour, draw=wpfour!70!black}]{WP4}{8}{16} \\
\ganttbar[bar/.append style={fill=wpfour!75, draw=wpfour!70!black}]{WP4.1}{8}{11} \\
\ganttbar[bar/.append style={fill=wpfour!75, draw=wpfour!70!black}]{WP4.2}{10}{14} \\
\ganttbar[bar/.append style={fill=wpfour!75, draw=wpfour!70!black}]{WP4.3}{12}{16} \\
\ganttgroup[group/.append style={fill=wpfive, draw=wpfive!70!black}]{WP5}{12}{18} \\
\ganttbar[bar/.append style={fill=wpfive!65, draw=wpfive!70!black}]{WP5.1}{12}{14} \\
\ganttbar[bar/.append style={fill=wpfive!65, draw=wpfive!70!black}]{WP5.2}{14}{16} \\
\ganttbar[bar/.append style={fill=wpfive!65, draw=wpfive!70!black}]{WP5.3}{16}{18}
\end{ganttchart}
\end{center}
}

\begin{center}
\small
\begin{tabular}{>{\raggedright\arraybackslash}p{0.13\textwidth} >{\raggedright\arraybackslash}p{0.82\textwidth}}
\toprule
\textbf{Código} & \textbf{Descripción resumida} \\
\midrule
WP1.1--WP1.3 & Cámara ambiental, integración de cámara rápida y flujo de adquisición abierto. \\
WP2.1--WP2.3 & Segmentación HSVM, CBF por FFT y CBP por flujo óptico/kymographs. \\
WP3.1--WP3.3 & Curaduría TEM, segmentación axonémica y cuantificación ODA/IDA. \\
WP4.1--WP4.3 & Curaduría genómica, priorización de variantes e integración fenotípica. \\
WP5.1--WP5.3 & Capacitación clínica, SOPs y liberación de software/informe final. \\
\bottomrule
\end{tabular}
\end{center}

\subsection*{6. Factibilidad, riesgos y mitigación}

\textbf{Factibilidad:} el proyecto usa infraestructura existente de FCEN-UBA (cómputo, taller, formación de recursos humanos) y colaboración clínica ya activa con Hospital Gutiérrez, reduciendo riesgo de implementación.

\textbf{Riesgos críticos y mitigación:}
\begin{itemize}[leftmargin=*]
\item \textit{Demoras de equipamiento:} priorización de compra local y plan de contingencia con hardware de laboratorio.
\item \textit{Heterogeneidad de datos:} estandarización de formatos y control de calidad automático en pipelines.
\item \textit{Variantes de significado incierto:} consenso ACMG/AMP con integración de evidencia funcional, estructural y clínica.
\end{itemize}

\subsection*{7. Impacto esperado y transferencia}

\begin{enumerate}[leftmargin=*]
\item Plataforma computacional transferida a hospital pediátrico público con trazabilidad diagnóstica.
\item Reducción de subjetividad analítica y mejora de reproducibilidad entre operadores.
\item Formación interdisciplinaria de estudiantes y profesionales clínicos.
\item Disponibilización regional en acceso abierto de software, SOPs y pipeline genómico \cite{Paff2021,Robson2025}.
\end{enumerate}

\subsection*{8. Referencias bibliográficas}

\begin{thebibliography}{9}
\bibitem{Robson2025} Robson EA \textit{et al}. \textit{ERJ Open Research}. 2025. DOI: \href{https://doi.org/10.1183/23120541.00466-2025}{10.1183/23120541.00466-2025}.
\bibitem{Shapiro2015} Shapiro AJ \textit{et al}. \textit{Pediatric Pulmonology}. 2015. DOI: \href{https://doi.org/10.1002/ppul.23304}{10.1002/ppul.23304}.
\bibitem{Collison2025} Collison R \textit{et al}. \textit{Clinical Medicine}. 2025. DOI: \href{https://doi.org/10.1016/j.clinme.2025.100319}{10.1016/j.clinme.2025.100319}.
\bibitem{ComiteNeumonologia2020} Comité Nacional de Neumonología (SAP). \textit{Archivos Argentinos de Pediatría}. 2020. DOI: \href{https://doi.org/10.5546/aap.2020.S164}{10.5546/aap.2020.S164}.
\bibitem{Bricmont2021} Bricmont N \textit{et al}. \textit{Diagnostics}. 2021. DOI: \href{https://doi.org/10.3390/diagnostics11091700}{10.3390/diagnostics11091700}.
\bibitem{Paff2021} Paff T \textit{et al}. \textit{Int J Mol Sci}. 2021. DOI: \href{https://doi.org/10.3390/ijms22189834}{10.3390/ijms22189834}.
\bibitem{Wrona2025} Wrona J \textit{et al}. \textit{J Clin Med}. 2025. DOI: \href{https://doi.org/10.3390/jcm14196808}{10.3390/jcm14196808}.
\bibitem{Balinotti2026} Balinotti JE \textit{et al}. \textit{Archivos Argentinos de Pediatría}. 2026. DOI: \href{https://doi.org/10.5546/aap.2025-10814}{10.5546/aap.2025-10814}.
\end{thebibliography}

\end{document}
